% Part: proof-theory
% Chapter: proof-search
% Section: introduction

\documentclass[../../../include/open-logic-section]{subfiles}

\begin{document}

\olfileid{pt}{ps}{int}

\olsection{مقدمه}

اصول و قواعد استنتاج دستگاه‌های !!{proof} مانند \Log{G3c} یا
\Log{N1i} تعیین می‌کنند چه درخت‌هایی از دنباله‌ها یا !!{formula}s
!!{proof}s درست به شمار می‌آیند. اگر چنین درختی از دنباله‌ها یا
!!{formula}s به ما داده شود---شاید همراه با اطلاعات افزوده، مانند
اینکه در هر گام چه قاعده‌ای اعمال شده، یا برچسب‌های فرض‌ها و قواعد
استنتاجی که آن‌ها را مرخص می‌کنند---تعریف اصول و قواعد به ما اجازه
می‌دهد بررسی کنیم آیا درخت داده‌شده !!{proof} درستی است یا نه. مسئله‌ای
متفاوت، و معمولاً دشوارتر، آن است که !!a{proof} برای دنباله‌ای معین، یا
برای !!{formula}ی معین از فرض‌هایی داده‌شده، \emph{بیابیم}.
این همان مسئلهٔ \emph{جست‌وجوی برهان} است.

روشی بسیار ساده‌انگارانه برای جست‌وجوی برهان وجود دارد. می‌توان
!!{formula}s را دنباله‌هایی از نمادهای یک الفبای متناهی دانست.
دنباله‌ها و درخت‌های دنباله‌ها یا !!{formula}s را نیز می‌توان
دنباله‌هایی از !!{formula}s دانست (شاید با پرانتزهای ویژه برای نمایش
شاخه‌شدن درخت). اصول و قواعد به ما امکان می‌دهند به‌طور مکانیکی بررسی
کنیم آیا دنباله‌ای از نمادها یک !!{proof} درست را بازنمایی می‌کند یا
نه. در نتیجه می‌توانیم \emph{همهٔ !!{proof}s درست ممکن} را به‌طور
مکانیکی تولید کنیم: همهٔ دنباله‌های نمادهای الفبای بازنمایی
!!{proof}s را تولید می‌کنیم و برای هر دنبالهٔ نامزد می‌سنجیم آیا
!!{proof} درستی را بازنمایی می‌کند. روش ساده‌انگارانهٔ جست‌وجوی برهان
چنین است: همهٔ !!{proof}s درست را تولید کنید تا یکی را بیابید که
!!{proof} دنبالهٔ داده‌شده باشد (یا برهانی برای !!{formula}
داده‌شده که فرض‌های بازش در میان فرض‌های داده‌شده باشند).

روش بهتر، استفاده از همان شیوهٔ ساده‌ای است که برای یافتن دستی
!!a{proof} به کار می‌برید: با دنبالهٔ پایانی آغاز می‌کنید،
!!a{formula}ی برمی‌گزینید و قاعده‌ای استنتاجی را «وارونه» اعمال
می‌کنید تا یک یا چند مقدمه تولید شود؛ مقدماتی که اگر !!{proof}s داشتند،
با آخرین استنتاج ترکیب می‌شدند و !!a{proof}ی برای دنبالهٔ داده‌شده
می‌ساختند. برای یافتن !!a{proof} در دستگاه استنتاج طبیعی نیز روشی
مشابه، هرچند پیچیده‌تر، به کار می‌رود.

اما این روش خطاناپذیر نیست: اگر قاعدهٔ نادرست یا !!{formula}s را به
ترتیب نادرست برگزینید، ممکن است به بن‌بست برسید. همچنین روش کاملاً
مشخص نشده است: در هر نقطه شاید بیش از یک !!{formula} ممکن برای انتخاب
یا بیش از یک قاعده برای اعمال وارونه وجود داشته باشد. \emph{الگوریتم
جست‌وجوی برهان} روشی کاملاً مشخص است که اگر !!a{proof}ی وجود داشته
باشد آن را می‌یابد؛ یعنی خطاناپذیر است.

برخی دستگاه‌های !!{proof} برای الگوریتم‌های ساده‌تر جست‌وجوی برهان
مناسب‌تر از دستگاه‌های دیگرند. در دستگاه‌های دنباله‌ای،
!!{main operator} یک !!a{formula} و اینکه در چپ یا راست رخ داده است،
قاعده‌ای را که باید وارونه اعمال شود تعیین می‌کند. برای نمونه، اگر
بخواهید !!a{proof}ی برای
$\Gamma \Sequent \Delta, !A \land !B$ بیابید که در آن
$!A \land !B$ در راست رخ می‌دهد، باید \RightR{\land} را وارونه اعمال
کنید و دو مقدمهٔ متناظر
$\Gamma \Sequent \Delta, !A$ و
$\Gamma \Sequent \Delta, !B$ را بسازید. اما اگر دستگاه شما
\Log{G1c} باشد، حضور انقباض کار را دشوارتر می‌کند، زیرا گزینه‌های
بیشتری می‌دهد: می‌توانید \RightR{\Contraction} را نیز وارونه اعمال
کنید و مقدمهٔ
$\Gamma \Sequent \Delta, !A \land !B, !A \land !B$ را بسازید؛ سپس
مقدمه‌ای با سه نسخه از~$!A \land !B$، و به همین ترتیب. در \Log{G2c}
وضع حتی بدتر است. اعمال وارونهٔ \RightR{\land} مستلزم حدس‌زدن
$\Gamma_1$ و $\Gamma_2$ چنان است که
$\Gamma = \Gamma_1 \cup \Gamma_2$، و $\Delta_1$ و $\Delta_2$ چنان که
$\Delta = \Delta_1 \cup \Delta_2$؛ سپس باید شانس خود را با مقدمات
$\Gamma_1 \Sequent \Delta_1, !A$ و
$\Gamma_2 \Sequent \Delta_2, !B$ بیازمایید. حدس نادرست ممکن است
بعدتر شما را به بن‌بست برساند؛ آن‌گاه باید به آغاز بازگردید و
$\Gamma_1$، $\Gamma_2$، $\Delta_1$ و $\Delta_2$ دیگری برگزینید. اگر
!!{formula} برابر $!A \lor !B$ باشد، باید یکی از دو صورت
\RightR{\lor} را برگزینید و یکی از دو مقدمهٔ ممکن
$\Gamma \Sequent \Delta, !A$ یا
$\Gamma \Sequent \Delta, !B$ را تولید کنید. انتخاب نادرست باز هم
شما را ناگزیر به پس‌گرد می‌کند.

دستگاه \Log{G3c} این مشکلات را ندارد. قاعدهٔ ساختاری ندارد و انتخاب
قاعدهٔ استنتاج با !!{main operator} و سمتی که !!{formula} در آن رخ
می‌دهد تعیین می‌شود. ازاین‌رو \Log{G3c} برای جست‌وجوی برهان مناسب‌تر
از \Log{G1c} یا \Log{G2c} است. جست‌وجوی موفق !!a{proof}ی در
\Log{G3c} به دست می‌دهد و سپس می‌توان آن !!{proof} را به \Log{G1c}
یا \Log{G2c} (و حتی \Log{N1c}) ترجمه کرد. پس داشتن الگوریتم جست‌وجوی
برهان برای \Log{G3c} و الگوریتم‌های ترجمهٔ !!{proof}s \Log{G3c} به
!!{proof}s دستگاه‌های دیگر، مسئلهٔ جست‌وجوی برهان را برای آن
دستگاه‌ها نیز حل می‌کند.

از کجا می‌دانیم الگوریتم جست‌وجوی برهان خطاناپذیر است؟ باید نشان دهیم
اگر الگوریتم نتواند !!a{proof}ی بیابد، هیچ !!{proof}ی نمی‌تواند وجود
داشته باشد. اگر الگوریتم بدی برگزینیم، ممکن است حتی با وجود برهان نیز
در یافتن آن شکست بخورد. الگوریتم ساده‌انگارانه چنین مشکلی ندارد، زیرا
همهٔ !!{proof}s ممکن را جست‌وجو می‌کند. اما الگوریتم جست‌وجوی برهان
باید کارآمد باشد و کمترین تلاش لازم را صرف کند.

ایدهٔ اصلی برای اثبات خطاناپذیری الگوریتم جست‌وجوی برهان این است که
نشان دهیم اگر الگوریتم شکست بخورد، از شیوهٔ شکست آن می‌توان نتیجه گرفت
که دنباله امکان ندارد !!a{proof}ی داشته باشد. برای نشان‌دادن این امر،
اثبات می‌کنیم دنباله معتبر نیست. در منطق مرتبهٔ اول کلاسیک، دنبالهٔ
$\Gamma \Sequent \Delta$ معتبر است اگر و تنها اگر هر !!{structure} دست‌کم
یک !!{formula} در $\Gamma$ را کاذب یا دست‌کم یک !!{formula} در
~$\Delta$ را صادق سازد. \Log{G3c} \emph{درست} است؛ یعنی تنها
دنباله‌های معتبر را !!{prove}s می‌کند. این مطلب با استقرا بر
!!{proof}s ثابت می‌شود: اصول آشکارا معتبرند و اگر مقدمات قاعده‌ای
معتبر باشند، نتیجهٔ آن نیز معتبر است. برای اثبات خطاناپذیری الگوریتم
نشان می‌دهیم اگر جست‌وجوی !!a{proof} برای
$\Gamma \Sequent \Delta$ شکست بخورد، می‌توانیم !!a{structure}ی تعریف
کنیم که در آن همهٔ !!{formula}s موجود در $\Gamma$ صادق و همهٔ
!!{formula}s موجود در~$\Delta$ کاذب‌اند. این کار همچنین \emph{تمامیت}
\Log{G3c} را اثبات می‌کند: اگر $\Gamma \Sequent \Delta$ معتبر باشد،
یک \Log{G3c}-!!{proof} برای آن وجود دارد. چون هر جست‌وجوی ناموفق نشان
می‌دهد $\Gamma \Sequent \Delta$ نامعتبر است، هر جست‌وجوی برهان برای
یک دنبالهٔ معتبر باید موفق شود.

\end{document}
